% Generated draft: requires mathematical and source review.
\documentclass[11pt]{article}
\usepackage[margin=25mm]{geometry}
\usepackage{fontspec}
\setmainfont{Latin Modern Roman}
\usepackage{amsmath,amssymb,mathtools,unicode-math}
\setmathfont{Latin Modern Math}
\usepackage{enumitem,needspace,xurl,hyperref}
\setlength{\parindent}{0pt}
\setlength{\parskip}{5pt}
\setlength{\emergencystretch}{3em}
\newcommand{\N}{\mathbb{N}}
\newcommand{\Z}{\mathbb{Z}}
\newcommand{\Q}{\mathbb{Q}}
\newcommand{\R}{\mathbb{R}}
\newcommand{\C}{\mathbb{C}}
\newcommand{\F}{\mathbb{F}}
\newcommand{\A}{\mathbb{A}}
\newcommand{\PP}{\mathbb{P}}
\newcommand{\E}{\mathbb{E}}
\newcommand{\eps}{\varepsilon}
\newcommand{\dd}{\,\mathrm{d}}
\begin{document}
\section*{UnsolvedMath ID: 30005125}
\textbf{Source problem code:} catalogue entry
\subsection{Definitions and mathematical statement}
The supplied record does not specify the ground field. The following definitions spell out the usual characteristic-zero conventions conditionally: let \(k\) be a field of characteristic zero. Whether the intended field is \(\C\), or another field with additional hypotheses, requires source clarification; the questions below are not asserted uniformly over all such fields.

The Virasoro Lie algebra \(\operatorname{Vir}\) has \(k\)-basis \(\{L_n:n\in\Z\}\cup\{z\}\), with
\[
[z,L_n]=0,\qquad [L_n,L_m]=(n-m)L_{n+m}+\delta_{n,-m}\frac{n^3-n}{12}z,
\]
where \(\delta_{n,-m}\) is the Kronecker delta. This fixes a conventional normalization, which is not specified in the supplied excerpt. Write \(S=S_k(\operatorname{Vir})\) for its symmetric algebra, equipped with the Poisson bracket determined by \(\{x,y\}=[x,y]\) on \(\operatorname{Vir}\) and the Leibniz rule. Write \(U=U_k(\operatorname{Vir})\) for its universal enveloping algebra.

An ideal \(I\subseteq S\) is a Poisson ideal if \(\{S,I\}\subseteq I\). For an ideal \(J\subseteq S\), its Poisson core \(\operatorname{PCore}(J)\) is the largest Poisson ideal contained in \(J\). A Poisson primitive ideal is an ideal of the form \(\operatorname{PCore}(\mathfrak m)\), where \(\mathfrak m\) is a maximal ideal of \(S\). Denote their set by \(\operatorname{PPrim}(S)\).

\begin{enumerate}
\item[(A)] Determine the inclusion relation on \(\operatorname{PPrim}(S)\): give a description or criterion deciding, for every pair \(P,Q\in\operatorname{PPrim}(S)\), whether \(P\subseteq Q\).
\item[(B)] For every nonzero scalar \(\lambda\) in the intended ground field, is the associative algebra
\[
U\big/\langle z-\lambda 1\rangle
\]
simple? Here \(\langle z-\lambda 1\rangle\) is the two-sided ideal generated by the central element \(z-\lambda 1\), and simple means nonzero with no two-sided ideals other than zero and the whole algebra.
\item[(C)] Does \(U\) satisfy the ascending chain condition on two-sided ideals? Explicitly, for every sequence of two-sided ideals
\[
I_1\subseteq I_2\subseteq I_3\subseteq\cdots\subseteq U,
\]
does there exist \(N\geq1\) such that \(I_n=I_N\) for every \(n\geq N\)?
\end{enumerate}

The supplied original excerpt additionally mentions a sub-question labelled (1a). For \(\chi\in\operatorname{Vir}^{*}=\operatorname{Hom}_k(\operatorname{Vir},k)\), define
\[
\mathfrak m_\chi=\langle x-\chi(x)1:x\in\operatorname{Vir}\rangle\subseteq S.
\]
That sub-question compares \(\operatorname{PCore}(\mathfrak m_\chi)\) with the principal ideal \((z-\chi(z)1)\), but its comparison symbol is corrupted as \texttt{\%}. Its intended relation, and the surrounding Question (1), cannot be recovered precisely from the supplied record.

As background, the excerpt reports a parameterization of Poisson prime ideals in \(S(W)=S(\operatorname{Vir})/(z)\), where \(W=\operatorname{Vir}/kz\) under the convention above, and attributes to its Theorem 4 the assertion that \(S(\operatorname{Vir})/(z-\lambda 1)\) is Poisson simple for \(\lambda\neq0\). A Poisson prime ideal is a prime ideal that is also Poisson; a Poisson algebra is Poisson simple if it is nonzero and has no nonzero proper Poisson ideals. These are supplied background claims, without the full theorem hypotheses or parameterization. They do not identify Poisson simplicity with the associative simplicity asked for in (B).
\subsection{Short English statement}
Describe which Poisson primitive ideals of the symmetric algebra of the Virasoro Lie algebra contain which others; ask whether every nonzero central reduction of its enveloping algebra is simple; and ask whether ascending chains of two-sided ideals in that enveloping algebra always stabilize. The intended ground field and a corrupted comparison in the ancillary question (1a) require clarification. The supplied partially solved label and background results are dated source assessments requiring human review.
\subsection{Sources}
\begin{itemize}
\item Local problem registry: supplied complete record for permanent UnsolvedMath ID \(30005125\), source code \texttt{catalogue entry}, titled \emph{Ideals and Central Reductions of the Virasoro Enveloping Algebra}; collection: Oberwolfach Reports Open Problems. No individual proposer is supplied. The proposed year is recorded as 2022.
\item Supplied report citation: \emph{Interactions between Algebraic Geometry and Noncommutative Algebra (2023)}, Oberwolfach Reports, pp.\ 1239--1312, DOI \texttt{10.4171/owr/2022/23}. The original excerpt identifies p.\ 1293. \url{https://doi.org/10.4171/owr/2022/23}. The citation year 2023 and the recorded proposed year 2022 are retained as supplied.
\item Supplied literature source, labelled ``Poisson spectrum paper'': \url{https://arxiv.org/abs/2106.02565}. The record supplies no author names or full bibliographic title for this link.
\item Supplied literature source, labelled ``Enveloping-algebra ideal paper'': \url{https://archive.intlpress.com/site/pub/files/_fulltext/journals/arkiv/2020/0058/0002/ARKIV-2020-0058-0002-a004.pdf}. The record supplies no author names or full bibliographic title for this link.
\item Literature assessment checked at \texttt{2026-08-22}: this is supplied dated curation, not independently verified. The supplied status is ``partially solved''. Its assessment says that the Poisson spectrum and Poisson simplicity of Virasoro central reductions are substantially understood, including Poisson simplicity for nonzero central parameter, while the stated associative simplicity and ascending-chain questions for enveloping-algebra quotients are not fully resolved. The supplied background describes this as dated literature triage derived from the OWR curation and requests independent verification of time-sensitive status and cited claims.
\item The registry labels the statement \texttt{corrected\_verified} and reports comparison with the official Oberwolfach report. This records the supplied verification claim; no independent source consultation or mathematical verification is asserted here.
\end{itemize}
\end{document}
